\documentclass[12pt]{article}
\usepackage[T1]{fontenc}
\usepackage[utf8]{inputenc}
\usepackage{lmodern}
\usepackage{textcomp}
\usepackage{enumitem}
\usepackage{geometry}
\geometry{margin=1in}
\begin{document}
\begin{center}
Answer Key - Philosophy - Graduate Level Assessment 11 \
\small (Answers are ordered exactly as in the question paper.)
\end{center}

\section*{Multiple Choice}
\begin{enumerate}[label=\arabic*.]
\item \textbf{Answer:} C
\\ \textit{Options:} A) art and beauty.; B) love and hatred.; C) virtue and vice.; D) success and failure.; E) truth and falsehood.
\\ \textit{Note:} Hume argues that our emotional responses to images of happiness and misery are essential for our moral sentiments, and without such emotional engagement, a person would lack the capacity to discern or care about the moral distinctions of virtue and vice.
\item \textbf{Answer:} D
\\ \textit{Options:} A) snob appeal; B) appeal to popularity; C) post hoc ergo propter hoc; D) appeal to anonymous authority; E) straw man
\\ \textit{Note:} The correct answer is "appeal to anonymous authority" because it refers to the fallacy of relying on the credibility of unnamed experts to support a claim without providing verifiable evidence or specific sources for their authority.
\item \textbf{Answer:} B
\\ \textit{Options:} A) the boundaries of our moral obligations.; B) the domain of our personal identity.; C) all of the above.; D) what we are free to do.; E) the responsibilities we have towards others.
\\ \textit{Note:} According to McGregor, rights are essential for defining and protecting the boundaries of personal identity, as they establish the moral and social frameworks within which individuals can express and realize their unique selves.
\item \textbf{Answer:} B
\\ \textit{Options:} A) it is impossible for humans to judge what is holy.; B) there are many other things that are also holy.; C) the concept of 'wrongdoing' is too subjective to define holiness.; D) only gods can decide who the wrongdoers are.; E) Euthyphro is not a reliable authority on what is holy.
\\ \textit{Note:} Socrates objects to Euthyphro's definition of the holy as prosecuting wrongdoers because it is too narrow and overlooks the existence of other actions and qualities that can also be considered holy, thereby failing to capture the essence of what is truly holy.
\item \textbf{Answer:} A
\\ \textit{Options:} A) to implement stricter guidelines for classroom discussions.; B) to promote greater understanding of historical and contemporary oppression.; C) to increase the number of safe spaces on campus.; D) to mandate sensitivity training for all students.; E) to increase funding for mental health services on campus.
\\ \textit{Note:} Lukianoff and Haidt argue that implementing stricter guidelines for classroom discussions can foster a more open environment for debate, thereby mitigating the effects of vindictive protectiveness that stifles intellectual discourse and encourages emotional responses over rational discussion.
\end{enumerate}

\section*{True / False}
\begin{enumerate}[label=\arabic*.]
\setcounter{enumi}{5}
\item \textbf{Answer:} False
\\ \textit{Options:} A) True; B) False
\\ \textit{Note:} Kamm argues that while enhancement can be considered a natural extension of medical treatment, it should not be prioritized over medical treatment, which is fundamentally aimed at restoring health and addressing medical needs.
\item \textbf{Answer:} True
\\ \textit{Options:} A) True; B) False
\\ \textit{Note:} The statement is true because Horyu-ji was indeed founded in 596 CE, making it the oldest surviving wooden temple in the world and a significant symbol of the introduction of Buddhism to Japan.
\item \textbf{Answer:} True
\\ \textit{Options:} A) True; B) False
\\ \textit{Note:} The retributive theory of punishment is based on the principle of just deserts, which asserts that punishment is justified as a moral response to wrongdoing rather than for its consequential benefits, thus it does not evaluate specific punishments based on the overall intrinsic value they produce.
\item \textbf{Answer:} True
\\ \textit{Options:} A) True; B) False
\\ \textit{Note:} Soto Zen emphasizes the practice of shikantaza, or "just sitting," in zazen meditation as a means to cultivate awareness and insight gradually, aligning with the concept of gradual enlightenment rather than sudden realization.
\item \textbf{Answer:} True
\\ \textit{Options:} A) True; B) False
\\ \textit{Note:} The statement is true because euthanasia is fundamentally understood as an intentional act aimed at ending a person's life to alleviate intractable suffering, reflecting ethical considerations surrounding autonomy and compassion in healthcare.
\end{enumerate}

\section*{Long Form Response}
\begin{enumerate}[label=\arabic*.]
\setcounter{enumi}{10}
\item \textbf{Answer:} Hobbes' concept of the state of nature posits a pre-political condition characterized by a "war of all against all," where individuals possess natural rights to do anything necessary for self-preservation. This chaotic environment necessitates the establishment of social contracts to ensure peace and security. Hobbes asserts that even in this state of war, every man retains a right to due process, suggesting that justice must prevail even amidst conflict. This notion implies that the legitimacy of power and authority is derived from a mutual agreement to uphold certain rights, including the right to a fair trial, thereby safeguarding individual dignity and fostering a semblance of order. Such an understanding challenges the absolutist tendencies in Hobbesian thought by emphasizing that the preservation of rights is fundamental to any governing structure, even in dire circumstances.
\\ \textit{Note:} correct because it accurately describes Hobbes' view of the state of nature as a chaotic environment where self-preservation is paramount, while also emphasizing his belief in the necessity of due process to maintain a sense of justice, even during a state of war.
\item \textbf{Answer:} Arthur argues that the rights of the poor to receive aid from the affluent cannot be deemed automatic or inherent, as they necessitate a contractual agreement. He posits that rights, particularly those involving redistributive justice, must be grounded in mutual consent and obligation, which implies that the affluent must explicitly agree to provide assistance. Without such a contract, the expectation of aid transforms into a moral claim rather than a binding right. Furthermore, this perspective emphasizes the importance of agency and responsibility in social relationships, suggesting that merely being affluent does not obligate individuals to support the poor unless they have entered into a clear, consensual agreement to do so. Thus, Arthur's view frames the discourse on poverty relief within a contractual context, delineating the boundaries of moral and legal responsibilities.
\\ \textit{Note:} correct because it captures Arthur's assertion that the rights of the poor to receive aid hinge on a contractual agreement grounded in mutual consent, distinguishing between moral claims and binding rights in the context of redistributive justice.
\end{enumerate}

\vfill
\noindent\textit{End of Answer Key}
\end{document}